\documentclass[11pt,a4paper]{article}

\usepackage[utf8]{inputenc}
\usepackage[T1]{fontenc}
\usepackage{amsmath,amssymb,amsfonts}
\usepackage{physics}
\usepackage{graphicx}
\usepackage{booktabs}
\usepackage{array}
\usepackage{siunitx}
\usepackage{hyperref}
\usepackage[margin=1in]{geometry}
\usepackage{xcolor}
\usepackage{listings}
\usepackage{enumitem}
\usepackage{fancyhdr}
\usepackage{titlesec}

\hypersetup{
    colorlinks=true,
    linkcolor=blue!60!black,
    urlcolor=blue!60!black,
    citecolor=blue!60!black,
    pdftitle={PY-BEMCS Physics Reference and User Manual},
    pdfauthor={Dr. Bharat Singh Rawat, Nick Magrin}
}

\pagestyle{fancy}
\fancyhf{}
\fancyhead[L]{PY-BEMCS Reference Manual}
\fancyhead[R]{\thepage}
\fancyfoot[C]{\small Python Beam Extraction \& Monte Carlo Simulator}

\titleformat{\section}{\Large\bfseries\color{blue!50!black}}{\thesection}{1em}{}
\titleformat{\subsection}{\large\bfseries\color{blue!40!black}}{\thesubsection}{1em}{}
\titleformat{\subsubsection}{\normalsize\bfseries\color{blue!30!black}}{\thesubsubsection}{1em}{}

\newcommand{\PYBEMCS}{PY-BEMCS}

\lstset{
    basicstyle=\small\ttfamily,
    backgroundcolor=\color{gray!10},
    frame=single,
    breaklines=true,
    captionpos=b
}

\begin{document}

% ---------------------- TITLE ----------------------
\begin{titlepage}
\centering
\vspace*{1.5cm}

{\Huge\bfseries Python Beam Extraction \&\\[0.3em] Monte Carlo Simulator \par}
\vspace{0.4cm}
{\Large\bfseries (\PYBEMCS)\par}

\vspace{1.0cm}
{\large\textit{Comprehensive Physics Reference, Numerical Controls, and User Manual}\par}

\vspace{1.8cm}
{\large\textbf{Original Author:}\par}
{\large Dr.~Bharat Singh Rawat\par}
{\normalsize \href{mailto:bharat.bharat22@gmail.com}{bharat.bharat22@gmail.com}\par}

\vspace{0.8cm}
{\large\textbf{Version 2.0 Extended Architecture \& Physics Modifications:}\par}
{\large Nick Magrin\par}
{\normalsize \href{mailto:nick.magrin@gmail.com}{nick.magrin@gmail.com}\par}

\vspace{1.2cm}
{\normalsize \textbf{Version 2.0 (Extended Modular Release)}\par}
\vspace{0.2cm}
{\normalsize \today\par}

\vfill
\hrule
\vspace{0.5em}
{\small Open-Source 2D3V Electrostatic Particle-in-Cell, Plume MCC, Thermo-Mechanical Grid Erosion, Numerical Verification Controls, \& Multi-Core Benchmarking Suite}
\end{titlepage}

% ---------------------- ABSTRACT ----------------------
\begin{center}
{\Large\bfseries Abstract}
\end{center}
\vspace{0.5em}

\PYBEMCS{} is an open-source 2D3V (two spatial, three velocity components) electrostatic Particle-in-Cell (PIC) and Monte Carlo Collision (MCC) simulation suite engineered for ion thruster optics, neutral beam injectors (NBI), and advanced plasma extraction systems. This document provides a complete mathematical, physical, and numerical specification of the simulator, together with rigorous stability criteria, physics validation checks, and comprehensive user guidelines for the graphical interface, headless batch solvers, and automated multi-core benchmarking scripts.

Key capabilities detailed in this manual include:
\begin{itemize}[nosep]
    \item Self-consistent Bohm extraction from upstream plasma sheaths with space-charge Poisson solves and Boltzmann electron screening.
    \item Multi-zone non-uniform finite-volume mesh with automatic geometry partitioning and adaptive axial coarsening factors ($1.0\times, 1.5\times, 2.0\times, 4.0\times$).
    \item Adaptive Picard iteration with dynamic residual convergence ($\Delta V_{\max} \le \SI{50}{\milli\volt}$) and sticky Newton--Raphson fallback solver.
    \item Automatic numerical physics validation controls (Debye length resolution, plasma frequency limits, CFL velocity bounds, and dynamic domain scaling).
    \item 2D3V relativistic Boris particle advancement with dual Taichi GPU/CPU accelerated and pure-NumPy fallback kernels.
    \item Automated parametric sweep framework with real-time multi-scenario switching and multi-metric steady-state convergence criteria.
    \item Charge-Exchange (CEX) collision dynamics using Roy's neutral plume model or user-imported experimental cross-section datasets with log-log cubic spline smoothing.
    \item In-situ sputter erosion modeling (Matsunami yield formulation) coupled with live grid mesh boundary deletion.
    \item Coupled 2D thermal conduction, Stefan--Boltzmann radiative cooling, and Euler--Bernoulli cantilever thermo-mechanical grid aperture deflection.
    \item Modular JSON configuration architecture (\texttt{config.json}) supporting arbitrary $N$-grid configurations, discharge chamber plasma dynamics, and live parameter reloading.
    \item Automated multi-core parallel benchmarking suite for CEX erosion rates, Electron Backstreaming (EBS) limits, primary beam impingement / grid scraping, and perveance-divergence curves.
    \item Built-in asynchronous standalone application packager (\texttt{PyInstallerWorker}) embedded directly in the GUI.
\end{itemize}

\vspace{1em}
\hrule
\vspace{1em}

\tableofcontents
\newpage

% =====================================================
\section{Computational Framework and Architecture}
% =====================================================

\subsection{Coordinate System and Multi-Zone Non-Uniform Discretisation}

The computational domain employs a 2D Cartesian grid in $(x, y)$, where $x$ represents the primary axial beam extraction coordinate and $y$ corresponds to the transverse/radial coordinate (half-aperture or axisymmetric slice). 

To optimize computational throughput while strictly resolving local physical scales (such as the sub-millimeter upstream Debye length $\lambda_{\mathrm{D}}$ and fine aperture grid chamfers), \PYBEMCS{} implements a \textbf{multi-zone non-uniform axial mesh} coupled with a uniform transverse spacing:
\begin{itemize}[nosep]
    \item Transverse spacing remains strictly uniform: $\Delta y = \text{const}$ across all $j \in [0, n_y - 1]$, where $n_y = \lfloor L_y / \Delta y \rfloor + 1$.
    \item Axial spacing $\Delta x(x)$ adapts across automatically detected physical zones:
    \begin{enumerate}[nosep]
        \item \textbf{Presheath Zone} ($x \in [0, x_{\mathrm{screen\_start}}]$): Baseline fine resolution $f_{\mathrm{pre}} = 1.0\times$ ($\Delta x = \Delta x_0 \le \lambda_{\mathrm{D}}$) to accurately capture Bohm meniscus formation.
        \item \textbf{Optics Grid Barrels} ($x \in [x_k, x_k + t_k]$): Moderate resolution $f_{\mathrm{barrel}} = 1.0\text{--}1.5\times$ ensuring $\ge 15$ cells across grid thickness $t_k$ for high-fidelity boundary tracking and sputter erosion.
        \item \textbf{Inter-Grid Acceleration Gaps} ($x \in [x_k + t_k, x_{k+1}]$): $f_{\mathrm{gap}} = 1.0\text{--}2.0\times$ resolving strong electrostatic acceleration fields while maintaining electric field gradient jumps below $10\%$.
        \item \textbf{Near and Far Plume Zones} ($x \in [x_{\mathrm{exit}}, L_x]$): Coarsened resolution $f_{\mathrm{plume}} = 2.0\text{--}4.0\times$, exploiting the dilute quasi-neutral state of the downstream beam to achieve a $25\text{--}35\%$ reduction in total grid nodes.
    \end{enumerate}
\end{itemize}

For each axial node $i$, the non-uniform control volume width $h_i^x$ is defined by:
\begin{equation}
    h_i^x = \frac{\Delta x_{i-1/2} + \Delta x_{i+1/2}}{2}, \qquad \Delta x_{i+1/2} = x_{i+1} - x_i.
\end{equation}
The 2D nodal cell volume is rigorously given by $\Omega_{i,j} = h_i^x \Delta y \Delta z$ (with unit transverse depth $\Delta z = \SI{1}{\milli\meter}$), ensuring exact machine-precision charge conservation ($|\sum Q_p - \iint \rho\, \dd{V}| / \sum Q_p < 10^{-8}$).

Every macroparticle stores three-dimensional velocity components $\vb{v} = (v_x, v_y, v_z)$. While spatial motion is constrained to the $(x, y)$ plane, $v_z$ is rigorously advanced in time during magnetic rotations and collisional scattering.

\subsection{Dual-Kernel Execution Engine (Taichi GPU + CPU Fallback)}

To combine peak GPU performance on workstation hardware with maximum portability across containerized or headless computing clusters, \PYBEMCS{} implements a dual-kernel architecture:
\begin{enumerate}[nosep]
    \item \textbf{Taichi Backend (GPU/CPU)}: When Taichi is initialized, particle pushes, bilinear charge deposition, and explicit thermal conduction compile into parallel GPU kernels (CUDA, Metal, or Vulkan) running in 32-bit floating-point precision ($\texttt{ti.f32}$).
    \item \textbf{Vectorised CPU Fallback Kernels}: If Taichi is not present (or when running inside frozen executable packages), the physics engine automatically activates vectorised CPU routines (\texttt{accumulate\_rho\_cpu}, \texttt{push\_particles\_boris\_cpu}) using optimized NumPy array operations.
\end{enumerate}

% =====================================================
\section{Rigorous Numerical Physics Controls and Stability Verification}
% =====================================================

To prevent unphysical numerical heating, artificial grid-aliasing instabilities, or phase-space skipping, the physics engine enforces strict automated validation checks during domain generation (\texttt{build\_domain}):

\subsection{Debye Length Resolution Criterion}
To prevent unphysical finite-grid self-heating, the spatial cell dimensions must resolve the upstream plasma Debye length:
\begin{equation}
    \lambda_{\mathrm{D}} = \sqrt{\frac{\varepsilon_0 T_{\mathrm{e,up}}}{e \cdot (0.61 n_0)}},
\end{equation}
where the pre-sheath factor $0.61 \approx \exp(-1/2)$ accounts for the density drop at the sheath edge. The engine strictly requires:
\begin{equation}
    \Delta x \le \lambda_{\mathrm{D}} \quad \text{and} \quad \Delta y \le \lambda_{\mathrm{D}}.
\end{equation}
If the cell spacing exceeds $\lambda_{\mathrm{D}}$, a descriptive error is raised, prompting the user to refine $\Delta x, \Delta y$ or adjust the plasma density.

\subsection{Plasma Frequency and Time-Step Integration Limits}
The time step $\Delta t$ must resolve both the ion plasma frequency $\omega_{\mathrm{pi}}$ and the scaled electron plasma frequency $\omega_{\mathrm{pe}}$:
\begin{equation}
    \omega_{\mathrm{pi}} = \sqrt{\frac{n_0 e^2}{m_{\mathrm{ion}} \varepsilon_0}}, \qquad
    \omega_{\mathrm{pe}} = \sqrt{\frac{n_0 e^2}{m_{\mathrm{e}} \varepsilon_0}}.
\end{equation}
The maximum permissible time step is bounded by:
\begin{equation}
    \Delta t \le \frac{2\pi}{\omega_{\mathrm{pi}}}, \qquad \Delta t_{\mathrm{e}} \le \frac{2\pi}{\omega_{\mathrm{pe}}}.
\end{equation}
If $\Delta t$ violates this limit, domain initialization halts to avoid numerical divergence.

\subsection{CFL and Velocity Resolution Bound}
Particles must not travel more than one mesh cell during a single time increment. The maximum expected particle velocity at the sheath edge includes the Bohm drift and 4 standard deviations of the thermal distribution:
\begin{equation}
    v_{\max} = v_{\mathrm{B}} + 4 v_{\mathrm{th,i}} = \sqrt{\frac{Z e T_{\mathrm{e}}}{m_{\mathrm{ion}}}} + 4\sqrt{\frac{Z e T_{\mathrm{i}}}{m_{\mathrm{ion}}}}.
\end{equation}
The engine enforces the Courant--Friedrichs--Lewy (CFL) velocity resolution bound:
\begin{equation}
    \frac{\Delta x}{\Delta t} \ge v_{\max} \quad \text{and} \quad \frac{\Delta y}{\Delta t} \ge v_{\max}.
\end{equation}

\subsection{Dynamic Domain Sizing and Symmetry Boundary Conditions}
The axial computational domain length $L_x$ is dynamically calculated based on the total physical grid stack:
\begin{equation}
    L_x = \SI{1.5}{\milli\meter} + \sum_{k=1}^{N} \left(t_k + g_k\right) + \SI{3.0}{\milli\meter},
\end{equation}
ensuring adequate space upstream for meniscus formation and downstream for plume expansion.

The transverse domain height $L_y$ and boundary conditions adapt to the selected geometry layout:
\begin{itemize}[nosep]
    \item \textbf{\texttt{half\_hole}} (Default): Half-pitch axisymmetric representation with specular symmetry reflection at $y = 0$ ($L_y = 0.5 r_{\mathrm{s}} + 0.30 \cdot \text{pitch}$).
    \item \textbf{\texttt{one\_hole}}: Full single-aperture extraction domain ($L_y = 3.0 r_{\mathrm{s}}$) with periodic boundary conditions along $y$.
    \item \textbf{\texttt{two\_holes}}: Dual-aperture inter-hole coupling domain ($L_y = 3.0 r_{\mathrm{s}} + \text{pitch}$) with periodic boundaries in $y$.
\end{itemize}

\subsection{Multi-Metric Steady-State Stationarity Criterion}
\label{sec:steady_state_criterion}

To guarantee that simulations and parametric sweep scenarios reach true physical steady state before recording telemetry or advancing to subsequent design points, the physics engine enforces a multi-metric stationarity evaluation across a rolling iteration window $W$ (default: $W = 100$ steps). 

Evaluating Poisson convergence or net charge balance alone is insufficient: sheath expansion and virtual anode relaxation can induce slow hydrodynamic current oscillations downstream long after the electrostatic potential appears stable. Therefore, stationarity is evaluated simultaneously across four fundamental physical observables:
\begin{enumerate}[nosep]
    \item \textbf{Total System Energy} ($E_{\mathrm{tot}}$): Sum of ion kinetic energy, electron kinetic energy, and electrostatic field energy:
    \begin{equation}
        E_{\mathrm{tot}} = \frac{1}{2} m_{\mathrm{ion}} w_{\mathrm{macro}} \sum_{p=1}^{N_{\mathrm{i}}} \|\vb{v}_p\|^2 + \frac{1}{2} m_{\mathrm{e}} w_{\mathrm{macro}} \sum_{e=1}^{N_{\mathrm{e}}} \|\vb{v}_e\|^2 + \frac{\varepsilon_0}{2} \int_{\Omega} \|\vb{E}\|^2 \dd{V}.
    \end{equation}
    \item \textbf{Total Net Space Charge} ($Q_{\mathrm{net}}$): Sum of discrete macroparticle charges and the interior Boltzmann electron fluid charge:
    \begin{equation}
        Q_{\mathrm{net}} = q_{\mathrm{ion}} w_{\mathrm{macro}} N_{\mathrm{i}} - e \, w_{\mathrm{macro}} N_{\mathrm{e}} - e \int_{\Omega_{\mathrm{plasma}}} n_0 \exp\!\left(\frac{V - V_{\mathrm{p}}}{T_{\mathrm{e}}}\right) \dd{V}.
    \end{equation}
    \item \textbf{Centerline Saddle-Point Potential} ($V_{\mathrm{saddle}}$): The minimum potential barrier along the aperture midline ($y = 0$) governing electron backstreaming and ion extraction sheath curvature.
    \item \textbf{Exit Ion Extraction Current} ($I_{\mathrm{ion,exit}}$): Instantaneous and cumulative ion current collected across the downstream Faraday cup exit boundary plane:
    \begin{equation}
        I_{\mathrm{ion,exit}} = \frac{q_{\mathrm{ion}} w_{\mathrm{macro}} \sum_{\mathrm{exit}} 1}{\Delta t}.
    \end{equation}
\end{enumerate}

For each observable $Q \in \{E_{\mathrm{tot}}, Q_{\mathrm{net}}, V_{\mathrm{saddle}}, I_{\mathrm{ion,exit}}\}$, the rolling relative variation over the evaluation window $W$ is computed as:
\begin{equation}
    \Delta Q[k] = \frac{Q[k] - Q[k-W]}{|Q[k-W]| + \epsilon} \times 100\%,
\end{equation}
where $\epsilon = 10^{-12}$ prevents division-by-zero singularities. Stationarity is confirmed if and only if all tracked quantities simultaneously satisfy:
\begin{equation}
    \max\!\left(|\Delta E_{\mathrm{tot}}|, \, |\Delta Q_{\mathrm{net}}|, \, |\Delta V_{\mathrm{saddle}}|, \, |\Delta I_{\mathrm{ion,exit}}|\right) \le \varepsilon_{\mathrm{threshold}},
\end{equation}
where the default convergence tolerance is $\varepsilon_{\mathrm{threshold}} = 1.0\%$. 

Furthermore, to prevent false premature triggers during initial beam transit, the stationarity check enforces:
\begin{itemize}[nosep]
    \item $k_{\mathrm{transit}} \ge k_{\mathrm{min}}$: At least $k_{\mathrm{min}} = 500$ iterations must elapse \emph{after} the first primary ions traverse the screen grid into the inter-grid optics.
    \item $k \ge W$: The scenario history must span at least the full evaluation window $W = 100$.
    \item Safeguard iteration limit $k_{\max} = 5000$: If stationarity is not reached within $k_{\max}$, the scenario terminates with a warning flag (\texttt{steady\_state\_reached = False}) to prevent indefinite run stalls.
\end{itemize}

\subsection{Adaptive Plume Remeshing Formulation}
\label{sec:adaptive_plume_remeshing}

To balance computational efficiency and field fidelity, \PYBEMCS{} features an automatic periodic \textbf{Adaptive Plume Remeshing} engine (\texttt{adaptive\_remesh\_settings.enabled = true}, default interval $N_{\mathrm{remesh}} = 150$ iterations):
\begin{itemize}[nosep]
    \item \textbf{Presheath \& Optics Invariance}: Upstream presheath ($[0, x_{\mathrm{screen}}]$) and inter-grid optics barrels/gaps strictly preserve their baseline resolution ($\Delta x_0$, factor $= 1.0$), ensuring that sheath curvature and inter-grid electric fields $\mathbf{E}$ remain unperturbed.
    \item \textbf{Automatic Plume Cell Sizing}: Downstream of the final grid exit ($x > x_{\mathrm{exit}}$), the local average ion density $n_{\mathrm{plume}}$ is sampled dynamically:
    \begin{equation}
        n_{\mathrm{plume}} = \frac{N_{\mathrm{plume,ions}} \cdot w_{\mathrm{macro}}}{V_{\mathrm{plume}}} \quad [\si{\per\meter\cubed}].
    \end{equation}
    The local Debye length expansion ratio determines the far-plume coarsening factor $f_{\mathrm{far}}$:
    \begin{equation}
        f_{\mathrm{target}} = \sqrt{\frac{n_0}{\max(n_{\mathrm{plume}}, \, 10^{12})}}, \quad f_{\mathrm{far}} = \operatorname{clip}\left(\operatorname{round}(f_{\mathrm{target}}, 1), \, 1.0, \, 8.0\right).
    \end{equation}
    The near-plume buffer zone smoothly transitions with factor $f_{\mathrm{near}} = \operatorname{clip}\left(\operatorname{round}(0.5(1.0 + f_{\mathrm{far}}), 1), \, 1.0, \, 2.0\right)$.
    \item \textbf{Dynamic Bidirectional Adaptation}: As beam density increases or focuses, the mesh cells automatically restrict (refine); as the plume dilutes or expands, cells automatically enlarge (coarsen).
    \item \textbf{Telemetry \& Visualization}: Tracked live in the Computational Domain \& Mesh Grid window via Tab 3 (\emph{Plume Cell Size History}) displaying $\Delta x_{\mathrm{plume}}$ vs.\ iterations alongside baseline nominal $\Delta x_0$.
\end{itemize}

\subsection{Child--Langmuir Perveance and Perveance Fraction Evaluation}
\label{sec:perveance_formulation}

Perveance characterizes the space-charge limited current extraction capacity of gridded ion optics. In \PYBEMCS{}, perveance metrics are evaluated automatically in all parametric sweeps, headless executions, and telemetry tables:

\begin{enumerate}[nosep]
    \item \textbf{Total Acceleration Voltage} ($V_{\mathrm{tot}}$):
    The net potential difference driving ion acceleration between the screen grid ($V_{\mathrm{s}}$) and accelerator grid ($V_{\mathrm{a}}$):
    \begin{equation}
        V_{\mathrm{tot}} = \max(V_{\mathrm{s}} - V_{\mathrm{a}}, \, 0.0).
    \end{equation}
    
    \item \textbf{Measured Effective Perveance} ($P_{\mathrm{eff}}$):
    Calculated directly from the measured steady-state exit ion beam current $I_{\mathrm{ion,exit}}$ and total acceleration voltage:
    \begin{equation}
        P_{\mathrm{eff}} = \frac{I_{\mathrm{ion,exit}}}{V_{\mathrm{tot}}^{3/2}} \quad \left[\si{\ampere\per\volt^{1.5}}\right].
    \end{equation}
    If $V_{\mathrm{tot}} \le 0$ or $I_{\mathrm{ion,exit}} \le 0$, $P_{\mathrm{eff}} = 0.0$.
    
    \item \textbf{Theoretical Maximum Child--Langmuir Perveance} ($P_{\max}$):
    Following Goebel and Katz~\cite{goebel}, the space-charge limited current limit for circular aperture optics accounts for screen grid thickness $t_{\mathrm{s}}$, inter-grid separation gap $l_{\mathrm{g}}$, and screen hole radius $r_{\mathrm{s}}$ through the equivalent acceleration length $l_{\mathrm{e}}$:
    \begin{equation}
        l_{\mathrm{e}}^2 = (l_{\mathrm{g}} + t_{\mathrm{s}})^2 + r_{\mathrm{s}}^2.
    \end{equation}
    The maximum theoretical Child--Langmuir perveance limit is:
    \begin{equation}
        P_{\max} = \frac{4\varepsilon_0}{9} \sqrt{\frac{2 q_{\mathrm{ion}}}{m_{\mathrm{ion}}}} \, \frac{A_{\mathrm{aperture}}}{l_{\mathrm{e}}^2} \quad \left[\si{\ampere\per\volt^{1.5}}\right],
    \end{equation}
    where $A_{\mathrm{aperture}}$ is the physical extraction open area of the screen grid hole:
    \begin{equation}
        A_{\mathrm{aperture}} =
        \begin{cases}
            \pi r_{\mathrm{s}}^2, & \text{in Axisymmetric ($r\text{--}z$) cylindrical coordinates}, \\
            r_{\mathrm{s}} \cdot d_{\mathrm{unit}}, & \text{in Cartesian 2D slab mode (\texttt{half\_hole}) with } d_{\mathrm{unit}} = \SI{1}{\milli\meter}.
        \end{cases}
    \end{equation}
    Critically, $A_{\mathrm{aperture}}$ is strictly defined by the aperture radius $r_{\mathrm{s}}$, remaining invariant to whether upstream macroparticle injection covers only the aperture or extends across the entire simulation domain ($R_{\mathrm{domain}} = L_y$). This avoids artificial $P_{\max}$ inflation and guarantees consistent perveance fractions between 2D Cartesian and radial simulations.
    
    \item \textbf{Perveance Fraction} ($\Pi$):
    The non-dimensional ratio quantifying beam extraction efficiency relative to the theoretical space-charge limit:
    \begin{equation}
        \Pi = \frac{P_{\mathrm{eff}}}{P_{\max}}.
    \end{equation}
    A perveance fraction $\Pi \ll 1$ indicates under-perveance (concave meniscus, excessive beam crossover and impingement), $\Pi \approx 0.2\text{--}0.6$ corresponds to optimal beamlet collimation with minimum divergence, and $\Pi \to 1.0$ approaches space-charge choking and virtual anode formation.
\end{enumerate}

% =====================================================
\section{Particle-in-Cell Electrostatic Algorithm}
% =====================================================

\subsection{Self-Consistent Poisson--Boltzmann Solver}

The electrostatic potential $\phi(\vb{x})$ is governed by Poisson's equation including macroparticle ion density and thermalized Boltzmann electrons:
\begin{equation}
    \nabla^{2}\phi(\vb{x}) \;=\; -\,\frac{1}{\varepsilon_{0}}
        \left[ q_{\mathrm{i}}\,n_{\mathrm{i}}(\vb{x}) - e\,n_{\mathrm{e}}(\vb{x}) \right],
    \label{eq:poisson}
\end{equation}
where $q_{\mathrm{i}} = Z e$ is the ion charge, $n_{\mathrm{i}}$ is the deposited ion macroparticle number density, and the electron density $n_{\mathrm{e}}$ follows the non-linear Boltzmann relation:
\begin{equation}
    n_{\mathrm{e}}(\vb{x}) \;=\; n_{0}\,
    \exp\!\left[\frac{\min(\phi(\vb{x}),\,\phi_{\mathrm{pl}}) - \phi_{\mathrm{pl}}}{T_{\mathrm{e}}}\right],
    \label{eq:boltzmann}
\end{equation}
with $T_{\mathrm{e}}$ in electron-volts ($\si{\eV}$), $n_0$ the upstream plasma density, and $\phi_{\mathrm{pl}}$ the reference plasma potential established at the upstream boundary ($x = 0$). The clipping operator $\min(\phi, \phi_{\mathrm{pl}})$ prevents non-physical exponential density blow-ups in regions where local potential excursions exceed the upstream plasma potential.

Equation~\eqref{eq:poisson} is discretised on the non-uniform mesh using a finite-volume formulation. Multiplying by $\Delta y^2$ yields the scaled 5-point stencil:
\begin{equation}
    \frac{\Delta y^2}{h_i^x} \left( \frac{\phi_{i+1,j} - \phi_{i,j}}{\Delta x_{i+1/2}} - \frac{\phi_{i,j} - \phi_{i-1,j}}{\Delta x_{i-1/2}} \right) +
    \left(\phi_{i,j+1} - 2\phi_{i,j} + \phi_{i,j-1}\right)
    = -\frac{\Delta y^2 \rho_{i,j}}{\varepsilon_{0}}.
\end{equation}
Scaling by $\Delta y^2$ preserves identical unit off-diagonals $(+1.0)$ along the $y$-direction, maintaining optimal matrix conditioning and sparse structure. The resulting sparse operator $\mathbf{A} \boldsymbol{\phi} = \mathbf{b}$ is pre-factorised using SuperLU sparse LU decomposition (or GPU-accelerated CuPy sparse LU).

\subsubsection{Adaptive Picard Iteration and Sticky Newton--Raphson Fallback}

To guarantee numerical stability and eliminate unphysical electric field spikes, the non-linear Boltzmann electron response is solved adaptively:
\begin{enumerate}[nosep]
    \item \textbf{Adaptive Picard Under-Relaxation}: At iteration $k$, the right-hand side space charge vector $\mathbf{b}^{(k)}$ incorporates the non-linear electron fluid density $\rho_e(\boldsymbol{\phi}^{(k)})$. The intermediate solution $\boldsymbol{\phi}^* = \mathbf{A}^{-1} \mathbf{b}^{(k)}$ is relaxed via:
    \begin{equation}
        \boldsymbol{\phi}^{(k+1)} = (1 - \omega) \boldsymbol{\phi}^{(k)} + \omega \boldsymbol{\phi}^*, \qquad \omega = 0.2.
    \end{equation}
    Iterations proceed dynamically until the maximum potential residual satisfies:
    \begin{equation}
        \|\Delta \boldsymbol{\phi}\|_{\infty} = \max_{i,j} \left|\phi_{i,j}^{(k+1)} - \phi_{i,j}^{(k)}\right| \le \SI{50}{\milli\volt}.
    \end{equation}
    \item \textbf{Sticky Newton--Raphson Fallback}: If Picard iteration stagnates (residual decreases by less than $1\%$ over 5 iterations) or exceeds 35 iterations, the engine automatically switches to a linearized Newton--Raphson solver:
    \begin{equation}
        \mathbf{J}(\boldsymbol{\phi}^{(k)}) \, \delta \boldsymbol{\phi} = -\mathbf{F}(\boldsymbol{\phi}^{(k)}), \qquad \boldsymbol{\phi}^{(k+1)} = \boldsymbol{\phi}^{(k)} + \delta \boldsymbol{\phi},
    \end{equation}
    where the Jacobian diagonal contains the Boltzmann derivative $J_{m,m} = A_{m,m} + \frac{\Delta y^2 e n_e(\phi_m)}{\varepsilon_0 T_e}$. Once engaged, the solver remains in Newton mode ('sticky') until steady-state convergence is re-established, achieving quadratic convergence without divergence risk.
\end{enumerate}

\subsection{Bilinear Charge Deposition}

Macroparticle charge is mapped to the four surrounding mesh nodes using bilinear interpolation weights. For a particle located at $(x_p, y_p)$ with grid indices $i = \lfloor x_p/\Delta x \rfloor$, $j = \lfloor y_p/\Delta y \rfloor$:
\begin{align}
    h_x &= \frac{x_p - i\Delta x}{\Delta x}, \qquad
    h_y = \frac{y_p - j\Delta y}{\Delta y},\\
    \rho_{i,j}     &\mathrel{+}= \frac{q_{\mathrm{i}} w}{V_{\mathrm{cell}}}\,(1-h_x)(1-h_y),\\
    \rho_{i+1,j}   &\mathrel{+}= \frac{q_{\mathrm{i}} w}{V_{\mathrm{cell}}}\,h_x(1-h_y),\\
    \rho_{i,j+1}   &\mathrel{+}= \frac{q_{\mathrm{i}} w}{V_{\mathrm{cell}}}\,(1-h_x)h_y,\\
    \rho_{i+1,j+1} &\mathrel{+}= \frac{q_{\mathrm{i}} w}{V_{\mathrm{cell}}}\,h_x h_y,
\end{align}
where $V_{\mathrm{cell}} = \Delta x \Delta y \Delta z$ is the cell volume with assumed unit transverse thickness $\Delta z = \SI{1}{\milli\meter}$, and $w$ is the statistical macroparticle weighting factor:
\begin{equation}
    w = \max\!\left(\frac{n_0 V_{\mathrm{cell}}}{N_{\mathrm{ppc}}},\; 10^{3}\right),
\end{equation}
with default nominal target $N_{\mathrm{ppc}} = 40$ particles per cell.

\subsection{Electric Field Computation and Boris Pusher}

Electric field components are evaluated at mesh nodes via second-order central differencing:
\begin{equation}
    E_{x,i,j} = -\frac{\phi_{i+1,j} - \phi_{i-1,j}}{2\Delta x}, \qquad
    E_{y,i,j} = -\frac{\phi_{i,j+1} - \phi_{i,j-1}}{2\Delta y},
\end{equation}
and interpolated to particle positions. Particle trajectories are advanced using the explicit leap-frog Boris method~\cite{birdsall}:
\begin{align}
    \vb{v}^{-} &= \vb{v}^{n-1/2} + \frac{q\,\vb{E}}{2m}\,\Delta t,\\
    \vb{t} &= \frac{q\,\vb{B}}{2m}\,\Delta t, \qquad
    \vb{s} = \frac{2\,\vb{t}}{1 + |\vb{t}|^2},\\
    \vb{v}' &= \vb{v}^{-} + \vb{v}^{-} \times \vb{t},\\
    \vb{v}^{+} &= \vb{v}^{-} + \vb{v}' \times \vb{s},\\
    \vb{v}^{n+1/2} &= \vb{v}^{+} + \frac{q\,\vb{E}}{2m}\,\Delta t,\\
    \vb{x}^{n+1} &= \vb{x}^{n} + \vb{v}^{n+1/2}\,\Delta t.
\end{align}
This scheme guarantees exact conservation of phase-space volume and energy conservation in static magnetic fields.

% =====================================================
\section{Plasma Injection and Bohm Optics}
% =====================================================

\subsection{Bohm Sheath Extraction Criterion}

Ions enter the computational domain through the upstream plasma boundary ($x = 0$) with an initial axial drift satisfying the Bohm sheath velocity criterion:
\begin{equation}
    v_{\mathrm{B}} = \sqrt{\frac{Z e T_{\mathrm{e}}}{m_{\mathrm{ion}}}}, \qquad
    m_{\mathrm{ion}} = M_{\mathrm{amu}} \cdot u,
    \label{eq:bohm}
\end{equation}
where $M_{\mathrm{amu}}$ is the ion mass in atomic mass units and $u = \SI{1.66054e-27}{\kilo\gram}$.

The injected ion current across the open sheath area $A_{\mathrm{inj}}$ is:
\begin{equation}
    I_{\mathrm{ion}} = 0.61\, Z e n_{0} v_{\mathrm{B}} A_{\mathrm{inj}},
\end{equation}
where the prefactor $0.61 \approx \exp(-1/2)$ accounts for the quasi-neutral pre-sheath density drop. The number of macroparticles injected at each timestep $\Delta t$ is:
\begin{equation}
    N_{\mathrm{inj}} = \frac{I_{\mathrm{ion}} \Delta t}{Z e w} + \xi_{\mathrm{stochastic}},
\end{equation}
with stochastic remainder accumulation ensuring exact long-term current conservation.

\subsection{Thermal Velocity Sampling}

Thermal spread at the sheath boundary is sampled from a 3D Maxwellian distribution at ion temperature $T_{\mathrm{i}}$:
\begin{equation}
    v_x = v_{\mathrm{B}} + \zeta_x v_{\mathrm{th,i}}, \qquad
    v_y = \zeta_y v_{\mathrm{th,i}}, \qquad
    v_z = \zeta_z v_{\mathrm{th,i}},
\end{equation}
where $\zeta_{x,y,z} \sim \mathcal{N}(0, 1)$ are standard normal random variables, and $v_{\mathrm{th,i}} = \sqrt{Z e T_{\mathrm{i}} / m_{\mathrm{ion}}}$.

% =====================================================
\section{Multiphysics Grid Models}
% =====================================================

\subsection{Arbitrary $N$-Grid Geometry Specification}

The extraction optics consist of an arbitrary sequence of $N$ conducting grids defined by:
\begin{itemize}[nosep]
    \item DC potential $V_k$ (\si{\volt})
    \item Thickness $t_k$ (\si{\milli\meter})
    \item Inter-grid axial gap to subsequent grid $g_k$ (\si{\milli\meter})
    \item Aperture radius $r_k$ (\si{\milli\meter})
    \item Upstream chamfer angle $\theta_k$ (\si{\degree})
\end{itemize}

Solid boundary masks $\mathcal{M}_k(x,y)$ are dynamically synthesised and rasterised onto the finite-difference Poisson mesh. However, physical boundary interactions, continuous sputtering erosion, and grid heat transfer are decoupled into dedicated continuous sub-grid models.

\subsection{Decoupled Sub-Grid Structural Surface Representation}
\label{sec:subgrid_surface}

In conventional PIC-MCC gridded optics codes, rasterizing curved apertures onto discrete Eulerian cells ($\Delta x \times \Delta y \sim \SI{20}{\micro\meter}$--\SI{50}{\micro\meter}) causes unphysical staircased corners that distort radial electric fields, over-scatter periphery ions, and result in chunky ``cell-dropping'' erosion. 

\PYBEMCS{} resolves this limitation by completely decoupling the physical grid geometry from the background Poisson mesh via a continuous sub-grid surface representation (\texttt{SubgridApertureProfile} and \texttt{SubgridSegment}):
\begin{itemize}[nosep]
    \item \textbf{Continuous Boundary Segments}: Each grid aperture is discretized into fine 1D line elements (\SI{5}{\micro\meter}--\SI{10}{\micro\meter} resolution) spanning:
    \begin{enumerate}[label=(\roman*),nosep]
        \item \emph{Upstream Face}: $x = g_{\mathrm{start}} + \delta_{\mathrm{cant}}$, $y \in [r_{\mathrm{front}}, L_y]$ with outward unit normal $\hat{\vb{n}} = (-1, 0)^{\mathsf{T}}$.
        \item \emph{Conical Chamfer / Barrel}: Connecting $(g_{\mathrm{start}} + \delta_{\mathrm{cant}}, r_{\mathrm{front}})$ to $(g_{\mathrm{end}} + \delta_{\mathrm{cant}}, r_{\mathrm{back}})$ with conical taper angle $\theta_{\mathrm{cham}}$ and inward-pointing normal:
        \begin{equation}
            \hat{\vb{n}} = (-\sin\theta_{\mathrm{cham}}, \, -\cos\theta_{\mathrm{cham}})^{\mathsf{T}}.
        \end{equation}
        \item \emph{Downstream Face}: $x = g_{\mathrm{end}} + \delta_{\mathrm{cant}}$, $y \in [r_{\mathrm{back}}, L_y]$ with outward normal $\hat{\vb{n}} = (+1, 0)^{\mathsf{T}}$.
    \end{enumerate}
    \item \textbf{Continuous Impact Angle Determination}: When an ion strikes the grid at continuous coordinates $(x_i, y_i)$ with velocity $\vb{v}_i$, it is projected onto the nearest sub-grid segment $k$. The exact incident angle $\theta_i$ is computed:
    \begin{equation}
        \cos\theta_i = \max\!\left(0.1, \; \frac{-\vb{v}_i \cdot \hat{\vb{n}}_k}{\|\vb{v}_i\|}\right).
    \end{equation}
\end{itemize}

\subsection{1D Radial Structural Thermal Continuum and Cantilever Deflection}
\label{sec:subgrid_thermal}

Rather than solving 2D finite-difference heat diffusion on coarse vacuum/solid cells, \PYBEMCS{} models each grid membrane as a dedicated 1D radial finite-volume continuum $r \in [r_{\mathrm{aperture}}, L_y]$ (\texttt{SubgridThermalModel}):
\begin{equation}
    \rho_k c_{p,k} t_k \frac{\partial T}{\partial t} = \frac{1}{r} \frac{\partial}{\partial r}\left(k_k t_k r \frac{\partial T}{\partial r}\right) + \dot{q}_{\mathrm{dep}}(r) - 2 \varepsilon_k \sigma_{\mathrm{SB}} \left(T^4 - T_{\infty}^4\right),
\end{equation}
where $t_k$ is grid thickness, $k_k$ is thermal conductivity, $\rho_k$ is material density, $c_{p,k}$ is specific heat capacity, $\varepsilon_k$ is surface emissivity, $\sigma_{\mathrm{SB}} = \SI{5.67e-8}{\watt\per\meter\squared\per\kelvin\tothe{4}}$, and $T_{\infty} = \SI{300}{\kelvin}$.

The boundary conditions are:
\begin{itemize}[nosep]
    \item \textbf{Inner Aperture Edge} ($r = r_{\mathrm{aperture}}$): Adiabatic / radiation-only condition, $\left.\frac{\partial T}{\partial r}\right|_{r = r_{\mathrm{ap}}} = 0$.
    \item \textbf{Outer Chassis Edge} ($r = L_y$): Isothermal clamped boundary anchored to the mounting ring at $T(L_y) = \SI{300}{\kelvin}$.
\end{itemize}

Numerical stability of explicit radial diffusion is strictly bounded by the 1D Fourier criterion:
\begin{equation}
    \mathrm{Fo} = \frac{k_k}{\rho_k c_{p,k}} \frac{\Delta t_{\mathrm{eff}}}{\Delta r^2} \le 0.20,
\end{equation}
preventing numerical oscillations during accelerated thermal integration ($\Delta t_{\mathrm{eff}} = \Delta t \times 10^7$).

The resulting mean grid temperature rise $\Delta \bar{T}_k = \bar{T}_k - T_{\infty}$ drives Euler--Bernoulli cantilever bending deflection of the circular grid aperture segment:
\begin{equation}
    \delta_{\mathrm{cant},k} = \frac{\alpha_k \Delta \bar{T}_k L_{\mathrm{cant},k}^2}{2 t_k}, \qquad L_{\mathrm{cant},k} = L_y - r_k,
\end{equation}
where $\alpha_k$ is the thermal expansion coefficient. This deflection dynamically updates grid axial positions ($g_{\mathrm{start}} + \delta_{\mathrm{cant}}$), accounting for thermal gap closure. Whenever cumulative structural deformation exceeds $\SI{5}{\micro\meter}$, the Poisson sparse matrix operator is automatically re-factored.

% =====================================================
\section{Collisions and Sputter Erosion}
% =====================================================

\subsection{Charge-Exchange (CEX) Dynamics and Plume Neutral Model}

Neutral gas density distribution throughout the grid gap and downstream region follows Roy's analytical plume expansion model~\cite{roy}:
\begin{equation}
    n_{\mathrm{n}}(x,y) = n_{\mathrm{n},0} \, a \left[1 - \frac{1}{\sqrt{1 + (r_T / R)^2}}\right] \cos\theta,
\end{equation}
with $R = \sqrt{y^2 + (x + r_T)^2}$, $\theta = \arctan(y / (x + r_T))$, and normalisation factor $a = (1 - 1/\sqrt{2})^{-1}$.

The null-collision Monte Carlo collision (MCC) probability for an ion traversing distance $\Delta s = |\vb{v}|\Delta t$ is:
\begin{equation}
    P_{\mathrm{CEX}} = 1 - \exp\left[-n_{\mathrm{n}}(x,y) \, \sigma_{\mathrm{CEX}}(E) \, |\vb{v}| \, \Delta t\right].
\end{equation}
The analytical xenon CEX cross-section formulation is:
\begin{equation}
    \sigma_{\mathrm{CEX}}(g) = \left[-0.8821\ln(g) + 15.1262\right]^2 \times 10^{-20} \;\si{\meter\squared},
\end{equation}
where $g = |\vb{v}|$ is relative collision speed. If custom cross-section CSV tables are imported, this formula is dynamically superseded by cubic spline interpolation in $\log_{10}(E)$--$\log_{10}(\sigma)$ space.

\subsection{Continuous Yamamura Sputter Erosion vs.\ Cell Removal}

Erosion of grid boundaries under ion bombardment is governed by the energy-dependent Matsunami yield~\cite{matsunami} combined with Yamamura's angular enhancement formulation~\cite{yamamura}:
\begin{equation}
    Y(E_i, \theta_i) = Y_{\mathrm{coeff}} \cdot \max\!\left(0, \, E_i - E_{\mathrm{th}}\right) \cdot f(\theta_i),
\end{equation}
where $Y_{\mathrm{coeff}}$ is the material sputtering coefficient, $E_{\mathrm{th}}$ is the threshold displacement energy, and the angular enhancement function is bounded to avoid numerical divergence at grazing incidence:
\begin{equation}
    f(\theta_i) = \min\!\left(3.0, \; \frac{1}{\cos\theta_i}\right), \qquad \cos\theta_i = \max\!\left(0.1, \; \frac{-\vb{v}_i \cdot \hat{\vb{n}}_k}{\|\vb{v}_i\|}\right).
\end{equation}

In legacy implementations, sputtering damage $D_{i,j}$ was accumulated onto discrete Cartesian grid cells. When $D_{i,j} \ge D_{\mathrm{fail}}$, entire cells were removed from the solid mask. This ``cell-dropping'' approach generated jagged, $90^\circ$ stepped notches of size $\Delta x \times \Delta y$, causing unphysical local electric field distortions and staircased corner scattering.

In the current architecture, sputtering erosion is evaluated directly on the decoupled sub-grid segments (\texttt{SubgridSegment}). Cumulative recession depth evolves continuously:
\begin{equation}
    \Delta d_k = \frac{Y(E_i, \theta_i) \, w_{\mathrm{macro}}}{D_{\mathrm{thresh}}} \, \Delta y,
\end{equation}
where $D_{\mathrm{thresh}}$ is the normalized damage threshold and $\Delta y$ is the local transverse mesh scale. This continuous formulation produces smooth, sub-micrometer erosion groove profiles that are fed directly to telemetry and contour diagnostics (\texttt{get\_groove\_profile()}) without altering Poisson mesh connectivity or introducing unphysical discrete cell deletions.

% =====================================================
\section{Modular Configuration Architecture (\texttt{config.json})}
% =====================================================

The simulator parameters are managed through a unified, hierarchical JSON structure (\texttt{config.json}). A representative configuration is shown below:

\begin{lstlisting}[language=json,caption={Standard \texttt{config.json} structure}]
{
  "beam_species": {
    "mass_amu": 131.293,
    "charge_state": 1
  },
  "grid_material": {
    "preset": "Molybdenum"
  },
  "simulation": {
    "n0_plasma": 7.73e17,
    "Te_up": 8.5,
    "Ti": 0.034,
    "Tn": 400.0,
    "n0": 5.27e18,
    "Accel": 0.5,
    "Thresh": 10000.0,
    "sim_mode": "Both",
    "geometry": "half_hole"
  },
  "discharge_chamber": {
    "radius_m": 0.02,
    "length_m": 0.10,
    "n_cusp": 4,
    "pitch_mm": 3.0
  },
  "grids": [
    { "V": 1300.0, "t": 1.2, "gap": 0.9, "r": 1.3, "cham": 20.0 },
    { "V": -150.0, "t": 0.6, "gap": 0.5, "r": 0.8, "cham": 0.0 },
    { "V": 50.0,   "t": 0.5, "gap": 4.0, "r": 1.2, "cham": 0.0 }
  ],
  "cross_sections": {
    "cx_file": "",
    "see_file": "",
    "custom_file": "",
    "spline_smoothing": 0.0
  },
  "terminal_output": {
    "grid_temperatures": true,
    "beam_divergence": true,
    "saddle_point_potential": true,
    "mean_particle_energy": true,
    "iteration_time": true
  }
}
\end{lstlisting}

% =====================================================
\section{Graphical User Interface and Features}
% =====================================================

\subsection{Menu Bar and Advanced Controls}

\begin{itemize}
    \item \textbf{Settings $\rightarrow$ Advanced Parameters...}: Configures neutralizer injection coordinates $(x_{\mathrm{n}}, r_{\mathrm{n}})$, plasma offset potential, and electron mass scaling ratio $\mathcal{R} = m_{\mathrm{ion}}/m_{\mathrm{e}}$.
    \item \textbf{Settings $\rightarrow$ Open Config JSON...}: Interactive file browser to load external JSON configuration profiles.
    \item \textbf{Settings $\rightarrow$ Reload config.json}: Hot-reloads configuration files from disk without restarting the application.
    \item \textbf{Settings $\rightarrow$ Build .exe...}: Compiles the current Python source code into a standalone single-file binary executable using an asynchronous background thread with live milestone tracking.
    \item \textbf{Beam $\rightarrow$ Ion Species...}: Configures beam ion atomic mass (amu) and ionization charge state ($+1, +2, \dots$) with predefined presets (Xe, Kr, Ar, N$_2$, O$_2$, H$_2$, He, Hg, Cs).
    \item \textbf{Beam $\rightarrow$ Cross-Section Manager...}: Import external experimental CSV tables for CEX, SEE yield, or custom reaction channels with interactive log-log cubic spline fitting.
    \item \textbf{Materials $\rightarrow$ Grid Material...}: Selects material presets (Molybdenum, Steel SS316, Titanium, Graphite) or custom thermophysical properties ($k, \rho, c_p, \varepsilon, \alpha, E, Y_0, E_{\mathrm{th}}$).
\end{itemize}

\subsection{Comprehensive Diagnostic Windows and Telemetry Suite}

The simulator provides an extensive suite of standalone, interactive diagnostic windows equipped with native Matplotlib navigation toolbars (Pan, Zoom, Save) and live data streaming:
\begin{itemize}[nosep]
    \item \textbf{Beam Divergence Angle ($\theta_{95}$ \& Midline)}: Evaluated from the 95\textsuperscript{th} percentile and midline trajectory angles at configurable axial evaluation planes downstream of the optics:
    \begin{equation}
        \theta_{95} = Q_{0.95}\!\left[\,\left|\arctan\left(\frac{v_y}{v_x}\right)\right|\,\right].
    \end{equation}
    \item \textbf{2D Plasma Number Density Visualizer (\texttt{density\_window.py})}: Full 2D logarithmic color contours of ion and electron density maps coupled with 1D radial profiles cut along two independent downstream axial planes.
    \item \textbf{Steady-State Convergence Monitor (\texttt{steady\_state\_window.py})}: Continuous rolling window variance analysis tracking key parameters ($\theta_{95}$, $V_{\mathrm{saddle}}$, $I_{\mathrm{beam}}$, $E_{\mathrm{kin}}$). Steady state is confirmed when rolling relative fluctuations drop below $\le 1.0\%$ over a configurable evaluation duration.
    \item \textbf{Electron Backstreaming (EBS) Saddle Potential}: Continuous monitoring of centerline potential in the accelerator aperture to verify backstreaming suppression ($V_{\mathrm{sp}} < -5\,\si{\volt}$).
    \item \textbf{Mesh Discretization \& Sub-Step Inspector (\texttt{mesh\_window.py})}: Visualizes axial zone boundaries, cell size transitions $\Delta x(x)$, and the real-time distribution of Boris sub-steps required by high-velocity particles.
    \item \textbf{Energy \& Charge Budget Conservation (\texttt{energy\_window.py}, \texttt{charge\_window.py})}: Tracks instantaneous kinetic, potential, and field energy conservation, as well as macroparticle vs.\ Boltzmann electron charge neutrality.
    \item \textbf{Ion Energy Distribution Function (\texttt{iedf\_window.py})}: Live spectral histogram tracking particle kinetic energy distributions exiting the extraction system.
    \item \textbf{Cross-Section Manager (\texttt{cs\_viewer\_window.py})}: Log--log visualizer and cubic-spline interpolation inspector for imported reaction cross-section tables.
    \item \textbf{Asynchronous Data Exporters}: Background progress dialogs for exporting time-series simulation telemetry to CSV, particle phase-space archives, and animated GIF recordings.
\end{itemize}

% =====================================================
\section{Multi-Core Benchmarking Suite (\texttt{benchmarks/})}
% =====================================================

The \texttt{benchmarks/} directory contains dedicated multi-core physics sweep scripts utilising Python's \texttt{multiprocessing} engine to execute parametric studies across independent CPU worker processes simultaneously:

\begin{table}[h!]
\centering
\small
\begin{tabular}{lp{6cm}p{4.5cm}}
\toprule
\textbf{Benchmark Script} & \textbf{Physics Objective} & \textbf{Generated Output} \\
\midrule
\texttt{benchmark\_cex.py} & Sputter erosion rate vs background neutral density ($n_{\mathrm{n}} = 10^{18}$--$5\times 10^{19}\,\si{\per\meter\cubed}$) across multiple grid gap configurations. & \texttt{cex\_benchmark\_multicore.png} \\
\addlinespace
\texttt{benchmark\_ebs.py} & Centerline saddle-point potential barrier vs accelerator grid voltage ($V_{\mathrm{a}} = -400$--$0\,\si{\volt}$) evaluating the EBS safety boundary. & \texttt{ebs\_benchmark\_multicore.png} \\
\addlinespace
\texttt{benchmark\_impingement.py} & Direct primary beam ion grid interception and scraping percentage vs perveance ($P = I_{\mathrm{ion}} / V_{\mathrm{tot}}^{1.5}$). & \texttt{impingement\_ratio\_benchmark.png} \\
\addlinespace
\texttt{benchmark\_perveance.py} & Beam divergence angle vs perveance over 60 plasma density increments across multiple grid gap spacings. & \texttt{divergence\_vs\_perveance\_benchmark\_multicore.png} \\
\addlinespace
\texttt{benchmark\_perveance\_Vs\_Sweep.py} & Beam divergence response as a function of screen grid extraction potential ($V_{\mathrm{s}} = 600$--$2000\,\si{\volt}$). & \texttt{divergence\_vs\_voltage\_benchmark.png} \\
\bottomrule
\end{tabular}
\caption{Summary of multi-core benchmarking scripts.}
\end{table}

% =====================================================
\section{Standalone Executable Compilation Architecture}
% =====================================================

The embedded executable builder (\texttt{Settings $\rightarrow$ Build .exe...}) runs PyInstaller asynchronously via the \texttt{PyInstallerWorker(QThread)} class. The compilation process proceeds without blocking the Qt event loop:
\begin{enumerate}[nosep]
    \item Spawns a dedicated background worker thread and builds inside a temporary directory (\texttt{\_exe\_build\_tmp/}) to keep the repository root clean.
    \item Intercepts real-time standard output and maps build milestones (analysis, dependency collection, PYZ archive creation, binary packaging) to a 0--100\% progress dialog.
    \item Automatically copies the resulting single-file executable and accompanying \texttt{config.json} into the user-selected destination directory.
    \item Cleans up intermediate build artifacts upon completion or cancellation.
\end{enumerate}

% =====================================================
\section{Summary of Execution Modes}
% =====================================================

\subsection{Interactive Graphical Run}
\begin{lstlisting}[language=bash]
python Python/main.py
\end{lstlisting}

\subsection{Headless Parameter Sweep Execution}
\begin{lstlisting}[language=bash]
python Python/run_simulation_from_config.py
\end{lstlisting}

\subsection{Parallel Multi-Core Benchmarks}
\begin{lstlisting}[language=bash]
python Python/benchmarks/benchmark_perveance.py
python Python/benchmarks/benchmark_cex.py
python Python/benchmarks/benchmark_ebs.py
\end{lstlisting}

% =====================================================
\begin{thebibliography}{9}
\bibitem{birdsall}
C.~K. Birdsall and A.~B. Langdon,
\textit{Plasma Physics via Computer Simulation},
Institute of Physics Publishing, Bristol, 1991.

\bibitem{roy}
R.~I.~S. Roy, \textit{Numerical Simulation of Ion Thruster Plume Backflow}, Ph.D. Thesis, Massachusetts Institute of Technology, Cambridge, MA, 1995.

\bibitem{matsunami}
N.~Matsunami \textit{et al.},
``Energy dependence of the yields of ion-induced sputtering of monatomic solids,'' \textit{Atomic Data and Nuclear Data Tables}, vol.~31, no.~1, pp.~1--80, 1984.

\bibitem{goebel}
D.~M.~Goebel and I.~Katz,
\textit{Fundamentals of Electric Propulsion: Ion and Hall Thrusters},
John Wiley \& Sons, Hoboken, NJ, 2008.

\bibitem{yamamura}
Y.~Yamamura, Y.~Itikawa, and N.~Itoh,
``Angular dependence of sputtering yields of monatomic solids,''
\textit{Institute of Plasma Physics Report IPPJ-AM-26}, Nagoya University, 1983.
\end{thebibliography}

\end{document}
